\chapter{Options à zéro jour et flux de gamma des dealers}\label{ch:s2:zero-day-options-and-dealer-gamma-flows}

En 2025, les options sur le S\&P 500 dont l'échéance tombe le jour même de leur négociation représentaient 59~\% du volume total des options sur le SPX, selon le décompte de Cboe. Les dealers qui prennent la contrepartie
se couvrent dans la journée, et leur couverture peut pousser l'indice plus loin ou le retenir. Baltussen, Da, Lammers et Martens ont établi que le
rendement des 30 dernières minutes avant la clôture est prédit par le rendement du reste de la journée, sur plus de 60 marchés de contrats à terme entre 1974 et 2020. Ils l'ont rattaché
à la couverture gamma des teneurs de marché d'options et des fonds à effet de levier. La stratégie qui en découle est facile à écrire~: suivre le mouvement de la journée quand les dealers sont
short en gamma, le contrer quand ils sont long. La difficulté consiste à savoir lequel. Ce chapitre simule des options à zéro jour dont la position des clients est tirée au hasard chaque jour,
estime le gamma des dealers à partir des positions ouvertes, comme le font les analystes, et mesure ce que coûte cette estimation. Sur le marché synthétique, la stratégie atteint un ratio de Sharpe
de 1.68 lorsque la position des dealers est connue et de 0.40 lorsqu'elle est devinée. Après coûts, l'hypothèse perd de l'argent. La construction repose sur \texttt{firm.gammaflow}.

\section{Options à zéro jour et qui les négocie}

Une option à zéro jour (livre~1, chapitre~26) est une option le jour de son échéance. Son gamma est concentré près du strike et croît sans limite à mesure que la clôture
approche. Un mouvement de quelques points à 15~h~45 modifie son delta bien davantage que le même mouvement à 10~h. Celui qui est vendeur de ces options détient une
position dont la couverture change rapidement et dont la taille dépend du niveau de l'indice.

\begin{dated}{2026-09}{Volume des options à zéro jour}\label{dat:s2:zero-day-options-and-dealer-gamma-flows:volume}
Cboe a indiqué que les options à zéro jour sur le SPX ont atteint en moyenne un record de 2.3~millions de contrats par jour en 2025, soit 59~\% de l'ensemble du volume d'options sur le SPX. En août 2025, leur
part a atteint un record de 62.4~\%, soit environ 2.4~millions de contrats par jour, les traders particuliers représentant selon les estimations 53~\% de ce volume.
\end{dated}

Qui détient quelle position compte davantage que le volume. Si les clients achètent les options, les dealers sont short en gamma~: ils achètent l'indice après sa hausse et vendent
après sa baisse. Si les clients les vendent, les dealers sont long en gamma et s'opposent au mouvement. Ni, Pearson, Poteshman
et White ont mis en évidence ce canal sur les actions individuelles. Le rééquilibrage de la couverture des teneurs de marché d'options, sans information sous-jacente, a agi sur
la volatilité des rendements des actions et la probabilité de fortes variations de cours.

\section{Estimer le gamma des dealers}

Les positions ouvertes par strike sont publiques. Le côté de chaque contrat, et donc qui est long et qui est short, ne l'est pas. Un analyste qui veut connaître le gamma des dealers doit
le supposer.

\begin{definition}[Estimation de l'exposition gamma]\label{def:s2:zero-day-options-and-dealer-gamma-flows:gex}
Une \emph{estimation de l'exposition gamma}\index{estimation de l'exposition gamma} (gamma exposure estimate) est le gamma total des dealers, calculé à partir des positions ouvertes par strike et d'une hypothèse
sur le côté que détiennent les dealers, en supposant en général qu'ils sont long des calls que vendent les clients et short des puts que ceux-ci achètent~: $\sum_K (\text{OI call}_K -
\text{OI put}_K)\,\Gamma_K$ selon cette convention.
\end{definition}

\begin{definition}[Niveau de bascule du gamma]\label{def:s2:zero-day-options-and-dealer-gamma-flows:flip}
Le \emph{niveau de bascule du gamma}\index{niveau de bascule du gamma} (gamma flip level) est le niveau de l'indice auquel une estimation du gamma total des dealers change de signe~: au-dessus de ce niveau
les dealers sont estimés long en gamma, en dessous ils sont short.
\end{definition}

\texttt{firm.gammaflow} simule 2\,000 jours. L'indice ouvre à 100, et les calls et les puts à zéro jour sont cotés sur 61 strikes espacés de 0.1, de 97 à 103.
Les positions ouvertes sur les calls sont centrées au-dessus de l'ouverture et celles sur les puts en dessous. Chacune est ajustée par un facteur aléatoire quotidien, avec davantage sur les strikes ronds. Chaque
jour, la position nette des clients en calls et en puts est tirée comme une part des positions ouvertes, entre $-1$ (short) et 1 (long). En moyenne, les clients
sont short de 0.2 des calls et long de 0.3 des puts, avec un écart-type de 0.5 d'un jour à l'autre, et les dealers détiennent l'opposé. Les
options des dealers gagnent en delta à mesure que l'indice bouge. Les dealers négocient un dixième de ce delta non couvert toutes les cinq minutes et tout ce qui reste sur les 30 dernières minutes,
et chaque transaction déplace le cours de la barre suivante (le \cref{lst:s2:zero-day-options-and-dealer-gamma-flows:days}).

Dans la réalité simulée, les dealers sont short en gamma à l'ouverture pour 56.4~\% des jours. L'estimation conventionnelle (dealers long de chaque call, short de chaque put) conclut que 51.5~\% des
jours sont short. Elle concorde avec le signe réel pour 56.0~\% des jours, et sa corrélation avec le gamma réel vaut 0.23
(la \cref{fig:s2:zero-day-options-and-dealer-gamma-flows:estimate}). Connaître le côté moyen des clients ne fait guère mieux, avec 57.5~\% et 0.24, parce que
ce qui échappe à l'estimation, c'est le côté propre à la journée et non la moyenne. L'estimation conventionnelle situe un niveau de bascule entre 97 et 103 pour 96.0~\% des
200 premiers jours, avec une médiane de 100.28. La réalité n'en présente un que pour 36.5~\% d'entre eux~: la plupart du temps, le gamma réel des dealers garde le même signe sur
toute la plage.

\begin{omfigure}
\begin{tikzpicture}
\begin{axis}[width=9cm,height=6.5cm,xlabel={\small gamma réel des dealers à l'ouverture},ylabel={\small estimation conventionnelle},
  tick label style={font=\footnotesize},xmin=-25,xmax=32,ymin=-25,ymax=30,grid=major,grid style={gray!25},table/col sep=comma]
\addplot[only marks,mark=*,mark size=0.8pt,omThm,opacity=0.6] table[x=true,y=convention]{figdata/strategies-2/06-zero-day-options-and-dealer-gamma-flows/estimate.csv};
\addplot[black,thin,domain=-25:30] {0};
\addplot[black,thin] coordinates {(0,-25) (0,30)};
\end{axis}
\end{tikzpicture}

\omcaption{Le gamma total réel des dealers à l'ouverture, face à l'estimation conventionnelle issue des positions ouvertes (dealers long de chaque call, short de chaque put),
les 400 premiers des 2\,000 jours synthétiques, en unités du modèle. Les points des quadrants supérieur gauche et inférieur droit sont les jours dont l'estimation
rate le signe~: 44~\% de l'ensemble des jours. Données~: \texttt{s2\_gammaflow.days}.}\label{fig:s2:zero-day-options-and-dealer-gamma-flows:estimate}
\end{omfigure}

\section{Aimantation au strike et poussée}

\begin{definition}[Aimantation au strike]\label{def:s2:zero-day-options-and-dealer-gamma-flows:pinning}
L'\emph{aimantation au strike}\index{aimantation au strike} (pinning) est la tendance d'un sous-jacent à clôturer près d'un strike où les positions ouvertes sont importantes le jour de l'échéance des options, attribuée
à la couverture des dealers qui sont long en gamma sur ce strike~: ils vendent quand le cours franchit le strike à la hausse et achètent quand il le franchit à la baisse.
\end{definition}

La poussée est l'autre régime~: les dealers short en gamma amplifient les mouvements. La simulation le montre clairement. Les jours où les dealers étaient réellement short en gamma à
l'ouverture, le rendement logarithmique de la journée avait un écart-type de 1.48~\%. Sur les jours long en gamma, il était de 0.81~\%, contre 0.99~\% pour les mêmes chocs sans aucune
rétroaction de couverture.

L'aimantation est une autre affaire. La part des clôtures à moins de 0.05 d'un strike rond (un multiple de 0.5) atteint 21.4~\% les jours long en gamma, 20.4~\% les jours
short en gamma et 21.4~\% sans rétroaction. Le seul hasard donne 20~\%. Le modèle amortit les mouvements quand les dealers sont long en gamma, mais rien dans celui-ci n'attire le cours vers un
strike particulier~: ici les dealers couvrent une part de leur delta non couvert, ce qui s'oppose au mouvement quel que soit le niveau de l'indice. Une aimantation suppose un gamma
concentré sur un seul strike dans les dernières minutes et une couverture qui le suit de près. Avant de trader une aimantation, il faut vérifier que le modèle, ou le marché,
en produit une.

\section{Trader le flux de couverture}

La stratégie suit Baltussen et ses coauteurs~: à 15~h~30, on se place long si l'indice a monté depuis l'ouverture, short s'il a baissé, et on clôture à 16~h.
Son rendement, en points de base de l'indice, selon le gamma des dealers à l'ouverture, et une règle qui suit le mouvement les jours short en gamma et le
contre les jours long en gamma~:

\begin{center}
\small
\begin{tabular}{@{}lccccc@{}}
\toprule
2\,000 jours, pb par jour & short en gamma (t) & long en gamma (t) & règle & Sharpe & après 1 pb\\
\midrule
côté réel & 3.12 (3.0) & $-3.40$ ($-4.1$) & 3.24 & 1.68 & 1.16\\
estimation conventionnelle & 1.02 (1.0) & $-0.52$ ($-0.5$) & 0.78 & 0.40 & $-0.11$\\
côtés moyens & 1.08 (1.3) & $-1.77$ ($-1.4$) & 1.28 & 0.66 & 0.14\\
\bottomrule
\end{tabular}
\end{center}

\begin{omfigure}
\begin{tikzpicture}
\begin{axis}[width=9cm,height=5.4cm,ybar,bar width=14pt,ylabel={\small 30 dernières minutes (pb par jour)},tick label style={font=\footnotesize},
  xmin=0.5,xmax=3.5,ymin=-4,ymax=4,xtick={1,2,3},xticklabels={côté réel,convention,côtés moyens},grid=major,grid style={gray!25},
  table/col sep=comma,legend style={font=\scriptsize,at={(0.5,-0.2)},anchor=north},legend columns=2]
\addplot[fill=omThm!70,draw=omThm] table[x=x,y=short]{figdata/strategies-2/06-zero-day-options-and-dealer-gamma-flows/regimes.csv};
\addplot[fill=omDef!60,draw=omDef] table[x=x,y=long]{figdata/strategies-2/06-zero-day-options-and-dealer-gamma-flows/regimes.csv};
\legend{dealers short en gamma,dealers long en gamma}
\end{axis}
\end{tikzpicture}

\omcaption{Le trade de momentum intrajournalier (le signe du mouvement de la journée jusqu'à 15~h~30 multiplié par le rendement des 30 dernières minutes) sur les 2\,000 journées synthétiques,
réparti selon le gamma des dealers à l'ouverture~: le côté réel, l'estimation conventionnelle et une estimation avec les côtés moyens des clients. Données~:
\texttt{s2\_gammaflow.regimes}.}\label{fig:s2:zero-day-options-and-dealer-gamma-flows:regimes}
\end{omfigure}

Avec le côté réel, la séparation entre régimes est nette (\cref{fig:s2:zero-day-options-and-dealer-gamma-flows:regimes})~: momentum les jours de gamma court,
retournement les jours de gamma long, et une règle qui utilise les deux rapporte 3.24~points de base par jour. L'estimation conventionnelle n'en conserve qu'un quart, soit 0.78~point de base
par jour, un ratio de Sharpe de 0.40 avant coûts et $-0.11$ après un point de base par aller-retour. Sans aucun régime, suivre le mouvement de la journée
rapporte 0.27~point de base par jour, avec un ratio de Sharpe de 0.14. Le mécanisme est réel dans la simulation. L'essentiel de sa valeur tient à une information que les données publiques ne
contiennent pas.

\section{Fiches de stratégie}

\begin{strategyfile}{Gamma long sous le niveau de bascule}\label{strat:s2:zero-day-options-and-dealer-gamma-flows:flip}
\sfield{Qui paie, et pourquoi} Personne ne paie~; la position est une protection achetée là où la couverture des dealers est estimée amplifier les mouvements.
\sfield{Instruments et plateformes} Options à courte échéance sur indice ou options sur contrats à terme.
\sfield{Signal} L'indice sous un niveau de bascule estimé.
\sfield{Dimensionnement et exécution} Petit~; le signe de l'estimation est faux un grand nombre de jours.
\sfield{Coûts} Fourchettes des options~; décote temporelle quand rien ne se passe.
\sfield{Comment elle meurt} Un niveau de bascule qui est un artefact de la convention de côté~: dans la simulation du chapitre, la convention en trouve un 96~\% des jours et
la vérité 36.5~\%.
\sfield{Horizon, capacité, infrastructure} De l'intrajournalier à quelques jours~; positions ouvertes par strike.
\sfield{Backtest honnête} Utiliser les positions ouvertes telles qu'elles étaient connues avant la séance.
\sfield{Sources} Aucune performance vérifiée.
\end{strategyfile}

\begin{strategyfile}{Pin trade près des grands strikes}\label{strat:s2:zero-day-options-and-dealer-gamma-flows:pin}
\sfield{Qui paie, et pourquoi} Acheteurs d'options arrivant à échéance près d'un strike où les dealers sont longs en gamma, si l'indice s'y maintient.
\sfield{Instruments et plateformes} Straddles courts ou butterflies au strike qui porte la plus forte position ouverte, en fin de séance d'échéance.
\sfield{Signal} Position ouverte importante à un strike voisin et gamma des dealers estimé long à cet endroit.
\sfield{Dimensionnement et exécution} Petit~; la perte si l'indice s'échappe est rapide et importante.
\sfield{Coûts} Fourchettes sur les options à l'échéance.
\sfield{Comment elle meurt} Pas d'aimantation~: dans le modèle du chapitre, les clôtures tombent près des strikes ronds aussi souvent que le hasard le prévoit.
\sfield{Horizon, capacité, infrastructure} La dernière heure de la séance d'échéance.
\sfield{Backtest honnête} Comparer à la part des clôtures proches d'un strike par le seul hasard.
\sfield{Sources} Aucune performance vérifiée.
\end{strategyfile}

\begin{strategyfile}{Momentum intrajournalier de couverture des dealers}\label{strat:s2:zero-day-options-and-dealer-gamma-flows:momentum}
\sfield{Qui paie, et pourquoi} Dealers et fonds à effet de levier qui doivent trader dans le sens du mouvement de la journée près de la clôture pour rééquilibrer leurs couvertures.
\sfield{Instruments et plateformes} Contrats à terme sur indice pendant les 30 dernières minutes.
\sfield{Signal} Le signe du rendement de la journée jusqu'à 15~h~30, conditionné au gamma estimé des dealers.
\sfield{Dimensionnement et exécution} Entrer à 15~h~30, sortir à la clôture.
\sfield{Coûts} Deux transactions sur contrats à terme par jour.
\sfield{Comment elle meurt} Une estimation fausse du côté des dealers~; un encombrement (crowding) sur la même clôture.
\sfield{Horizon, capacité, infrastructure} Trente minutes~; données intrajournalières.
\sfield{Backtest honnête} Coûts sur chaque journée~; l'estimation de régime telle qu'elle aurait pu être faite ce matin-là.
\sfield{Sources} Baltussen, Da, Lammers et Martens (2021)~; ce chapitre~: 1.68 avec le côté réel, 0.40 avec la convention, $-0.11$ après coûts.
\end{strategyfile}

\section{Tutoriel~: qui est court en gamma}\label{tut:s2:zero-day-options-and-dealer-gamma-flows:tut}

\begin{tutorial}
\textbf{Objectif.} Simuler des options à zéro jour avec des côtés clients aléatoires, estimer le gamma des dealers à partir des positions ouvertes, et trader le flux de couverture
selon le régime. \textbf{État final~:} le tableau et les deux figures.
\begin{tutsteps}
\item \textbf{Les journées}.
\omcode{code/firm/gammaflow/firm_gammaflow.py}{87}{118}{Positions ouvertes, côtés clients, couverture des dealers avec retards, et l'estimation selon la convention.}\label{lst:s2:zero-day-options-and-dealer-gamma-flows:days}
\item \textbf{Le trade par régime}.
\omcode{code/strategies-2/06-zero-day-options-and-dealer-gamma-flows/python/s2_gammaflow.py}{49}{71}{Le trade des 30 dernières minutes, réparti selon le régime réel et le régime estimé.}\label{lst:s2:zero-day-options-and-dealer-gamma-flows:trade}
\item \textbf{Exécuter} \texttt{estimates()}, \texttt{regimes()}, \texttt{day\_vol()}, \texttt{pinning()}, \texttt{flips()} et
\texttt{fig\_gammaflow.py}.
\end{tutsteps}
\textbf{Pour aller plus loin.} Faire couvrir aux dealers le delta de chaque option à chaque barre et voir si l'aimantation apparaît~; donner à l'analyste un signal bruité
du côté client de la journée et mesurer le trade à mesure que le signal s'améliore~; ajouter le rééquilibrage de clôture des fonds à effet de levier.
\end{tutorial}

\section{Construction~: flux de gamma}\label{bld:s2:zero-day-options-and-dealer-gamma-flows:gammaflow}

\begin{build}
\bfield{Rôle} Positions ouvertes à zéro jour par strike, côtés clients, gamma des dealers et son estimation conventionnelle, niveaux de bascule, et un
simulateur intrajournalier dans lequel la couverture des dealers fait bouger les cours.
\bfield{Interface} \texttt{FlowConfig(\dots)}, \texttt{strikes(cfg)}, \texttt{open\_interest(cfg, centre)}, \texttt{gamma(S, K, tau, vol)},
\texttt{dealer\_gamma(S, tau, pos\_c, pos\_p, cfg)}, \texttt{flip\_level(pos\_c, pos\_p, cfg, tau)}, \texttt{simulate\_days(cfg)}.
\bfield{Règles} Les dealers détiennent l'opposé des clients~; les trades de couverture font bouger la barre suivante~; l'estimation ne voit que les positions ouvertes.
\bfield{Tests d'acceptation} \texttt{code/firm/gammaflow/tests/}~: le gamma culmine au strike et augmente près de l'échéance~; sans impact, on obtient les trajets libres~; un niveau de bascule
entre les calls longs et les puts courts.
\bfield{Pour aller plus loin} Couverture par option~; flux des fonds à effet de levier~; signal de côté bruité.
\end{build}

\begin{omsources}
\item Cboe, trading volume reports and Cboe Insights on SPX zero-day options, 2025.
\item G.~Baltussen, Z.~Da, S.~Lammers et M.~Martens, ``Hedging demand and market intraday momentum'', \emph{Journal of Financial Economics}
142(1), 2021.
\item S.~X.~Ni, N.~D.~Pearson, A.~M.~Poteshman et J.~White, ``Does option trading have a pervasive impact on underlying stock prices?'',
\emph{Review of Financial Studies} 34(4), 2021.
\end{omsources}

\section{Exercices}

\begin{exercise}[$\star$]\label{exo:s2:zero-day-options-and-dealer-gamma-flows:1}
Les dealers sont courts en gamma de 5 unités d'indice par point d'indice. L'indice monte de 2 points. Que doivent-ils trader pour rester couverts~?
\end{exercise}

\begin{exercise}[$\star$]\label{exo:s2:zero-day-options-and-dealer-gamma-flows:2}
Pourquoi le gamma d'une option à zéro jour près du strike augmente-t-il à mesure que la clôture approche~?
\end{exercise}

\begin{exercise}[$\star$]\label{exo:s2:zero-day-options-and-dealer-gamma-flows:3}
Le seul hasard place 20~\% des clôtures à moins de 0.05 d'un multiple de 0.5. Pourquoi~?
\end{exercise}

\begin{exercise}[$\star\star$]\label{exo:s2:zero-day-options-and-dealer-gamma-flows:4}
Pourquoi l'estimation conventionnelle trouve-t-elle un niveau de bascule presque chaque jour~?
\end{exercise}

\begin{exercise}[$\star\star$]\label{exo:s2:zero-day-options-and-dealer-gamma-flows:5}
Connaître les côtés moyens des clients n'améliore l'estimation que très peu. Pourquoi~?
\end{exercise}

\begin{exercise}[$\star\star$]\label{exo:s2:zero-day-options-and-dealer-gamma-flows:6}
Pourquoi la volatilité de la journée diffère-t-elle entre les jours à gamma court et les jours à gamma long, et que disent Ni et ses coauteurs à ce sujet~?
\end{exercise}

\begin{exercise}[$\star\star\star$]\label{exo:s2:zero-day-options-and-dealer-gamma-flows:7}
\emph{Programmation.} Régler \texttt{FlowConfig(impact=0.0)} et relancer la transaction. Que devrait-elle rapporter, et pourquoi~?
\end{exercise}

\begin{exercise}[$\star\star\star$]\label{exo:s2:zero-day-options-and-dealer-gamma-flows:8}
\emph{Trouver la faille.} «~Notre estimation du gamma indique que les dealers sont short aujourd'hui, donc les trente dernières minutes suivront le mouvement du matin.~»
\end{exercise}

\section{Problème~: qui est short gamma}

\begin{problem}[{Problème du week-end --- couvertures des dealers}]\label{pb:s2:zero-day-options-and-dealer-gamma-flows:1}
Les journées simulées du chapitre, les chiffres de Cboe et la recherche publique.

\medskip\noindent\textbf{Partie I --- Les options.}
\begin{enumerate}
\item Qu'est-ce qu'une option à zéro jour, et comment son gamma évolue-t-il au cours de la journée~?
\item Donner les chiffres de volume de Cboe pour 2025.
\item Comment les dealers à gamma court et à gamma long se couvrent-ils~?
\item Que trouvent Ni, Pearson, Poteshman et White~?
\end{enumerate}

\medskip\noindent\textbf{Partie II --- L'estimation.}
\begin{enumerate}[resume]
\item Définir une \omterm{def:s2:zero-day-options-and-dealer-gamma-flows:gex}{estimation de l'exposition gamma} et le niveau de bascule.
\item Quelle convention les analystes emploient-ils, et que suppose-t-elle~?
\item À quelle fréquence l'estimation trouve-t-elle le bon signe, et comment se corrèle-t-elle avec la réalité~?
\item À quelle fréquence chacune trouve-t-elle un niveau de bascule~?
\end{enumerate}

\medskip\noindent\textbf{Partie III --- Le marché.}
\begin{enumerate}[resume]
\item Comment les dealers se couvrent-ils dans le modèle, et comment leur couverture fait-elle bouger les prix~?
\item Donner la volatilité de la journée dans chaque régime.
\item Définir l'\omterm{def:s2:zero-day-options-and-dealer-gamma-flows:pinning}{aimantation au strike} (pinning). Le modèle produit-il cette aimantation~?
\item Que faudrait-il à un modèle pour produire une aimantation~?
\end{enumerate}

\medskip\noindent\textbf{Partie IV --- Le verdict.}
\begin{enumerate}[resume]
\item Énoncer le \emph{résultat nommé}~: l'erreur de l'estimation du gamma lorsque le côté des dealers est deviné, et le rendement de la transaction sur les flux de couverture par
régime.
\item Que trouvent Baltussen et ses coauteurs~?
\item Que rapporte la transaction sans distinction de régime~?
\item Qu'est-ce qui subsiste après les coûts~?
\item D'où viendraient de meilleures informations sur les côtés~?
\item Comment backtester honnêtement la transaction~?
\item Quelle fiche de stratégie dépend le plus du niveau de bascule~?
\item En une phrase~: sur quoi la transaction parie-t-elle réellement~?
\end{enumerate}
\end{problem}

\section{Questions d'entretien}

\begin{interviewq}[$\star$ \iqroles{trader}]\label{iq:s2:zero-day-options-and-dealer-gamma-flows:1}
Que signifie, pour les dealers, être short gamma, et pourquoi cela importerait-il pour l'indice~?
\end{interviewq}

\begin{interviewq}[$\star\star$ \iqroles{chercheur}]\label{iq:s2:zero-day-options-and-dealer-gamma-flows:2}
Comment estimeriez-vous le gamma des dealers, et comment sauriez-vous si l'estimation est bonne~?
\end{interviewq}

\begin{interviewq}[$\star\star$ \iqroles{trader}]\label{iq:s2:zero-day-options-and-dealer-gamma-flows:3}
L'indice se trouve près d'un strike où les positions ouvertes sont énormes une heure avant l'échéance. Que pourrait-il se passer, et comment opéreriez-vous~?
\end{interviewq}

\begin{interviewq}[$\star\star$ \iqroles{risque}]\label{iq:s2:zero-day-options-and-dealer-gamma-flows:4}
Un desk vend chaque jour des options à zéro jour. Quel est son risque, et comment le limitez-vous~?
\end{interviewq}

\begin{interviewq}[$\star\star$ \iqroles{développeur}]\label{iq:s2:zero-day-options-and-dealer-gamma-flows:5}
Concevoir un système qui calcule une \omterm{def:s2:zero-day-options-and-dealer-gamma-flows:gex}{estimation de l'exposition gamma} chaque minute à partir des positions ouvertes et des cotations.
\end{interviewq}

\begin{interviewq}[$\star\star\star$ \iqroles{chercheur}]\label{iq:s2:zero-day-options-and-dealer-gamma-flows:6}
Montrer que le gamma d'une option à la monnaie vaut $\varphi(\tfrac12\sigma\sqrt\tau)/(S\sigma\sqrt\tau)$, et déterminer comment il se comporte lorsque $\tau \to 0$.
\end{interviewq}
